\subsection{Contracting an internal flow edge}
The new operation changes the divergence energy rather than only the
alias penalties. It removes a flow coordinate on an internal tree edge
and identifies its endpoints. It is therefore not an alias identification
or an instance of square polarization. Its soundness is unconditional;
coefficient height one is a separate exact acceptance test.

\begin{theorem}\label{thm:contraction}
Let an incidence presentation have the energy form
\eqref{eq:networkenergy}, with every internal edge belonging to a tree
inside one row block or the hub block. Choose an internal edge from $u$
to $v$, with its own flow coordinate $w$. Delete that edge, identify its
vertices, and replace their two divergences by their sum. Retain all
prescribed occurrence labels and equality and unit penalties. Then the
old and new integer zero sets are in polynomial bijection, preserving
the original arithmetic coordinates. The new presentation has one fewer
flow coordinate and degree at most four. If full collection of its
polynomial has height one, exact degree four, and every listed variable
occurs, it is a certified height-one quartic on that smaller coordinate set.
\end{theorem}
\begin{proof}
Write the old two divergences as
\[
 d_u=a_u+\epsilon w,\qquad d_v=a_v-\epsilon w,
 \qquad \epsilon\in\{-1,1\},
\]
where $a_u,a_v$ are integer polynomials independent of $w$. No other
divergence or penalty depends on $w$, because it is the independent
flow coordinate of this one edge. The merged divergence is
$d_{uv}=a_u+a_v$. Thus the new polynomial is exactly
\[
 P'=\frac12\left((a_u+a_v)^2+
           \sum_{t\ne u,v}d_t^2\right)+\sum C+\sum Q+U.
\]
This expression proves nonnegativity and the degree bound without any
claim that the old and new polynomials themselves agree.

At an old zero, nonnegativity forces $d_u=d_v=0$ and every other
summand to vanish. Deleting $w$ therefore gives $a_u+a_v=0$ and a new
zero. Conversely, at a new zero all its displayed summands vanish.
Set
\[
 w=\epsilon a_v.
\]
Then $d_v=0$, while $d_u=a_u+a_v=0$; all other old summands are unchanged.
Indeed the stronger identity $P|_{w=\epsilon a_v}=P'$
holds for every remaining tuple: the two old divergences become
$a_u+a_v$ and zero. This is an integral polynomial restoration, and
it is unique on zeros because
$\epsilon$ is a unit. Restoration and deletion are mutually inverse on
the zero sets. All equality and unit arguments, and hence all recovered
logical coordinates, are identical on the two sides.

The edge's endpoints belong to one block. Summing the contracted
vertex divergences in that block therefore gives the same boundary sum
as before contraction. In particular \eqref{eq:networkboundary}
continues to recover exactly the same normalized source row; no source
occurrence is deleted. Contracting an edge of a tree leaves a tree.
The remaining flows are still uniquely recovered by leaf elimination,
or equivalently by composing the one-edge restoration above.

Finally, degrees and coefficient heights concern the polynomial after
all monomial collisions, including equality and unit terms. Contracting
vertices can create parallel prescribed incidences, so the initial simple
graph argument alone is no longer a sufficient height certificate. The
selection test reconstructs and collects the entire new energy and checks
its actual integer coefficients. This establishes the stated height,
exact degree and coordinate census only for the accepted finite instance.
\end{proof}

\begin{corollary}\label{cor:contractionsequence}
Any finite certified sequence of accepted internal-edge contractions
preserves the calibrated integer zero set. Its inverse is obtained by
restoring the deleted flows in reverse order. If all $k$ deleted flows
were distinct live coordinates and no other coordinate is removed, its
independent-variable count decreases by exactly $k$.
\end{corollary}
\begin{proof}
Compose the polynomial deletion and restoration maps of
Theorem~\ref{thm:contraction}. Every step preserves the block boundary
identity and leaves internal blocks as trees. The final independently
collected array fixes the actual coordinate count and all measured
polynomial invariants.
\end{proof}

This argument does not predict which edges pass the height test or
minimize the support. A contraction can save a variable while creating
more monomials; it can also fail the height test entirely. Adaptive
selection is a construction method, not a proof of optimality. The
released certificate freezes each accepted edge, so reproduction requires
no heuristic search.
